\chapter*{Notation}
\addcontentsline{toc}{chapter}{Notation}

\noindent
\begin{xltabular}{\textwidth}{@{}lX@{}}
\toprule
\textbf{Symbol} & \textbf{Meaning} \\
\midrule
\endfirsthead
\toprule
\textbf{Symbol} & \textbf{Meaning} \\
\midrule
\endhead
\bottomrule
\endfoot
\bottomrule
\endlastfoot
$n, N$ & Width of a hidden layer (number of neurons / channels); $N \in \hNinf$ denotes an unlimited hypernatural width. \\
$d, D$ & Depth, or number of hidden layers or blocks, excluding the input layer; $D \in \hNinf$ denotes an unlimited hypernatural depth. \\
$n_0, N_0$ & Input dimension (standard $n_0 \in \N$, or unlimited $N_0 \in \hNinf$ in Chapter~3). \\
$m, M$ & Number of inputs / training examples / tokens (standard $m \in \N$, or unlimited $M \in \hNinf$ in Chapter~3 and Section~6.2.4). \\
$x^\alpha$ & Input (data point) indexed by $\alpha \in \{1, \dots, m\}$. \\
$y^\alpha$ & Label / target associated with $x^\alpha$. \\
$\theta$ & Network parameters (weights). \\
$W_0, W_\ell, W_d$ & For scalar-output MLPs of hyperfinite width $N$, $W_0 \in \hR^{N\times n_0}$, $W_\ell \in \hR^{N\times N}$ for $1 \le \ell \le d-1$, and $W_d \in \hR^{1\times N}$ is the readout row vector. \\
$f(x^\alpha; \theta) = f^\alpha(\theta)$ & Network output at input $x^\alpha$. \\
$\varphi$ & Activation function (with natural extension ${}^*\varphi : \hR \to \hR$, often written $\varphi$ when unambiguous). \\
$c$ & Variance-preserving normalization constant, $c := (\E[\varphi(w)^2])^{-1}$ for $w \sim \mathcal{N}(0,1)$. \\
$h_\ell^\alpha$ & Hidden pre-activations (or hidden states) at layer $\ell$ for input $x^\alpha$. \\
$\varphi_\ell^\alpha = \varphi(h_\ell^\alpha)$, $\varphi_\ell = [\varphi_\ell^1 \;\cdots\; \varphi_\ell^m]$ & Hidden post-activations at layer $\ell$ and their dataset-column matrix. \\
$g_\ell^\alpha$ & Backward post-activation neurons at layer $\ell$ for input $x^\alpha$. \\
$z_\ell^\alpha$ & Backward pre-activation neurons at layer $\ell$ for input $x^\alpha$. \\
$\Phi^{(N)}, \Phi_\ell^{(N)}, \Phi_\ell^{\prime(N)}$ & Hyperfinite empirical feature and derivative kernels in Chapter~2, with hidden-layer normalization $1/N$ and input normalization $1/n_0$. The superscript $(N)$ records hyperfinite width. \\
$\Phi, \Phi_\ell, \Phi_\ell'$ & Standard-part kernels in Chapters~2 and~4 ($\Phi_\ell = \st(\Phi_\ell^{(N)})$). In Chapters~5 and~6, $\Phi_\ell$ also denotes internal layer covariances and $\Phi_\tau$ their standard SDE shadow. \\
$K_t^{(N)}, K_t, K$ & Hyperfinite empirical NTK and standard-part NTKs. Subscripts record time and superscripts in parentheses record width. \\
$L(\theta)$ & Training objective (empirical risk); $L(\theta) = \frac{1}{m}\sum_{\alpha=1}^m \ell(f^\alpha(\theta), y^\alpha)$. \\
$\E, \Prob$ & Standard expectation and probability. \\
$\norm{\cdot}, \ip{\cdot}{\cdot}$ & Euclidean norm and inner product (standard or internal by transfer, unless specified otherwise). \\
\midrule
$\hR, \hN, \hNinf$ & Robinson's $\aleph_1$-saturated hyperreals, hypernaturals, and unlimited hypernaturals. \\
$\Lim, \Inf, \App$ & Ring of limited hyperreals $\Lim$, maximal ideal of infinitesimals $\Inf$, and group of appreciable hyperreals $\App = \Lim \setminus \Inf$; $a \approx b \iff a - b \in \Inf$, and $a \asymp b \iff a/b \in \App$. \\
$\st, \mon(x)$ & Standard-part homomorphism $\st : \Lim \to \R$ ($\st(x) \in \R$ unique with $x \approx \st(x)$) and monad (halo) $\mon(x) = \{y \in \hR : y \approx x\}$. \\
$\hP, \hE, \Loeb, \EL$ & Internal probability $\hP$, internal expectation $\hE$, Loeb probability measure $\Loeb$ on $\mathcal{L}(\mathcal{A})$, and Loeb expectation $\EL$; $X \approxP Y$ means $X \approx Y$ $\Loeb$-almost surely. \\
\midrule
$\Cl_m, \hCl_N$ & Standard Euclidean Clifford algebra $\Cl_m = \Cl(\R^m)$ and internal hyperfinite Clifford algebra $\hCl_N = \Cl(\hR^N)$; for standard $m \in \N$, $\st : \Lim(\hCl_m) \to \Cl_m$ is the $\R$-algebra shadow homomorphism on limited multivectors. \\
$u v = \scalar{u v} + u \wedge v$, $\rev{M}$ & Geometric product of vectors decomposing into scalar inner part $\scalar{uv} = \ip{u}{v}$ and exterior bivector $u \wedge v$, with reversion $\rev{(v_1 \cdots v_k)} = v_k \cdots v_1$. \\
$\grade{M}{k}$, $A \lcontr B$, $A \rcontr B$ & Grade-$k$ projection $\grade{M}{k}$ and left/right interior contractions; for an invertible blade $\bm{B}$, $P_{\bm{B}}(x) = (x \lcontr \bm{B})\bm{B}^{-1}$ and $P_{\bm{B}}^\perp(x) = (x \wedge \bm{B})\bm{B}^{-1}$ are orthogonal projection and rejection. \\
$\dirac_\theta$, $\Spin(m)$, $\Pin(m)$ & Vector derivative (Dirac operator) $\dirac_\theta = \sum_{p} e_p \partial_{\theta_p}$ satisfying $\dirac_\theta^2 = \Delta_\theta$, and the even/full rotor covering groups $\Spin(m) \subset \Cl_m^+$, $\Pin(m) \subset \Cl_m^\times$. \\
\end{xltabular}
